\chapterrule
\section{反向传播究竟在传什么？}
\phantomsection\label{concept:backprop}
\label{intro:17-backpropagation}

多层网络的参数通过后续计算逐步影响最终损失。
若为每个参数都重新展开整条计算链，会做许多重复运算。
\textbf{反向传播}要解决的，是复用同一张计算图里的局部导数，
一次高效求出所有参数的梯度。

\subsection{从一条计算链理解链式法则}

先看一台三级机器：
\[
x\xrightarrow{\times\,w_1}h
\xrightarrow{\times\,w_2}\widehat{y}
\xrightarrow{\text{算损失}}L.
\]
令 $h=w_1x$、$\widehat{y}=w_2h$、$L=\tfrac12(\widehat{y}-y)^2$。
这里只有一个样本，所以不再像前几章那样对三点求平均；
系数 $\tfrac12$ 只为让导数 $\partial L/\partial\widehat{y}=\widehat{y}-y$ 更简洁。
现在要问：调整 $w_1$ 一点点，最后的损失 $L$ 会变多少？

最直接的做法是把整个函数展开，但神经网络可以有数十乃至更多层。链式法则沿着计算路径，把各步的\textbf{局部变化率}相乘：
\begin{equation}
\boxed{\frac{\partial L}{\partial w_1}
=\frac{\partial L}{\partial\widehat{y}}
\frac{\partial\widehat{y}}{\partial h}
\frac{\partial h}{\partial w_1}.}
\label{eq:chain-three}
\end{equation}
这叫\textbf{链式法则}（chain rule）。

\subsection{亲手算一次两层网络的梯度}
\label{sec:backprop-weight-gradients}

取 $x=2$、真实目标 $y=5$、初始权重 $w_1=1$、$w_2=2$。
前向过程：
\[
h=1\times2=2,\quad
\widehat{y}=2\times2=4,\quad
L=\tfrac12(4-5)^2=\tfrac12.
\]

从输出向输入反向算：
\[
\frac{\partial L}{\partial\widehat{y}}
=\widehat{y}-y=-1,\qquad
\frac{\partial\widehat{y}}{\partial h}=w_2=2,\qquad
\frac{\partial h}{\partial w_1}=x=2.
\]
所以
\[
\boxed{\frac{\partial L}{\partial w_1}
=(-1)\times2\times2=-4.}
\]
而对于最后一层权重：
\[
\frac{\partial L}{\partial w_2}
=\frac{\partial L}{\partial\widehat{y}}
\frac{\partial\widehat{y}}{\partial w_2}
=(-1)\times h=-2.
\]
两条梯度都为负：略增 $w_1$ 或 $w_2$，预测 $4$ 都会向目标 $5$ 靠近。
\hyperref[sec:end-to-end-step]{本章最后一节}给这个网络加上两项偏置，
只需补算偏置的梯度，再让四个参数一起更新。

\begin{center}
\begin{tikzpicture}[>=Latex,
  block/.style={draw=linegray,fill=paper,rounded corners=2pt,
  minimum width=34mm,minimum height=15mm,align=center,font=\small}]
\node[block] (h) at (0,0) {$h=w_1x$\\前向：$1\times2=2$};
\node[block] (out) at (4.25,0) {$\widehat y=w_2h$\\前向：$2\times2=4$};
\node[block] (loss) at (8.5,0) {$L=\frac12(\widehat y-5)^2$\\前向：$0.5$};
\draw[->,thick,blue] (h)--(out);
\draw[->,thick,blue] (out)--(loss);
\node[font=\footnotesize,teal,align=center] (g1) at (0,-1.45)
 {$\partial L/\partial h=(-1)\times2=-2$\\$\partial L/\partial w_1=(-2)\times2=-4$};
\node[font=\footnotesize,teal,align=center] (g2) at (4.25,-1.45)
 {$\partial L/\partial\widehat y=-1$\\$\partial L/\partial w_2=(-1)\times2=-2$};
\draw[->,teal,thick] (loss.south) |- (g2.east);
\draw[->,teal,thick] (g2.west) -- (g1.east);
\end{tikzpicture}

{\footnotesize 蓝箭头是前向值，青色箭头表示从损失回传梯度；
参数在梯度全部算完以后统一更新。}
\end{center}

\begin{idea}
反向传播（backpropagation）传的不是文字、答案或原始输入，而是损失对中间变量的梯度。
前向时从 $x$ 算到 $L$；求梯度时从 $\partial L/\partial\widehat{y}$ 出发，沿同一条链倒着走，这就是“反向”的含义。
后面的局部变化率先算出来，前面的参数直接复用，避免每求一个权重都从头展开整条函数。
\end{idea}\phantomsection\label{route:backprop-end}

\subsection{如果中间有 ReLU，怎样求导？}

先只看 ReLU 这一步：输入记作 $z$，输出记作 $h$，即
\[
h=\relu(z)=\max(0,z).
\]
求导就是问：\textbf{输入稍微变化时，输出跟着变化多少？}
图~\ref{fig:relu-function} 中的两段直线，可以分开来看。

\textbf{输入为正时，$h=z$。}
例如 $z$ 从 $2$ 增加到 $2.1$，$h$ 也从 $2$ 增加到 $2.1$。
输出变化量除以输入变化量是 $0.1/0.1=1$，所以这一段的导数为 $1$。

\textbf{输入为负时，$h=0$。}
例如 $z$ 从 $-2$ 增加到 $-1.9$，$h$ 始终为 $0$。
输出变化量除以输入变化量是 $0/0.1=0$，所以这一段的导数为 $0$。
只要输入的变化留在同一段内，把变化量缩得再小，这两个比例也不变。因此
\[
\frac{\partial h}{\partial z}=\relu'(z)=
\begin{cases}
0,&z<0,\\
1,&z>0.
\end{cases}
\]

\textbf{恰好在 $z=0$ 时，两段在这里形成折角。}
从左侧靠近，斜率是 $0$；从右侧靠近，斜率是 $1$。
两侧不一致，所以这一点没有通常意义下的导数。
但反向传播程序仍需要一个数来完成下面的乘法。
本书把零点按负数一侧处理，约定：\textbf{若输入恰好为 $0$，反传经过 ReLU 时也乘 $0$。}
这是折角处采用的计算规则，并不是求出了该点的导数。

现在把 ReLU 接回网络：
$z=w_1x+b_1$、$h=\relu(z)$、$\widehat{y}=w_2h+b_2$。
先从输出端算出 $\partial L/\partial h$；反传经过 ReLU 时，乘上刚求出的局部导数：
\[
\frac{\partial L}{\partial z}
=\frac{\partial L}{\partial h}\frac{\partial h}{\partial z},\qquad
\frac{\partial L}{\partial w_1}
=\frac{\partial L}{\partial z}\frac{\partial z}{\partial w_1}
=\frac{\partial L}{\partial z}\,x.
\]

沿用前面的 $x=2,y=5,w_1=1,w_2=2$，两项偏置取 $0$。
此时 $z=2,h=2,\widehat{y}=4$，所以
\[
\frac{\partial L}{\partial h}=(4-5)\times2=-2,\qquad
\frac{\partial L}{\partial w_1}=(-2)\times\underbrace{1}_{\text{ReLU 导数}}\times2=-4.
\]
只把输入改成 $x=-2$，目标和参数不变，则 $z=-2,h=0,\widehat{y}=0$，此时
\[
\frac{\partial L}{\partial h}=(0-5)\times2=-10,\qquad
\frac{\partial L}{\partial w_1}=(-10)\times\underbrace{0}_{\text{ReLU 导数}}\times(-2)=0.
\]
虽然预测仍有误差，乘积中的 ReLU 导数为 $0$，使这个样本经这条路径传给 $w_1$ 的梯度为 $0$。
这就是“梯度被截断”的含义：负数段的输出保持不变，前面参数的微小变化暂时无法沿这条路径影响损失。
其他不经过这个 ReLU 的路径仍可能提供梯度。

\subsection{反向传播负责求梯度，优化器负责更新}

\begin{center}
\tagpdfsetup{table/header-rows={1}}
\begin{tabularx}{\linewidth}{@{}p{29mm}X@{}}\toprule
概念 & 它负责什么\\\midrule
前向传播 & 给定参数，算输出和损失\\
反向传播 & 高效计算损失对每个参数的梯度\\
梯度下降/优化器 & 根据梯度修改参数\\
\bottomrule\end{tabularx}
\end{center}

更复杂的计算图会有分叉和相加：同一个参数影响损失的多条路径需要把各条路径的梯度贡献相加。自动微分框架正是在计算图上应用这个规则。

\begin{warningbox}
每个训练步都用同一组旧参数完成前向与反向传播，算出全部梯度，再由优化器统一更新所有参数。
不要反传到某一层就立刻修改那一层，否则同一步会混用新旧权重。
\end{warningbox}
